\documentclass[11pt,a4paper]{article}
\usepackage{fontspec,amsmath,amssymb,amsthm,microtype}
\setmainfont{lmroman10-regular.otf}[BoldFont=lmroman10-bold.otf,ItalicFont=lmroman10-italic.otf,BoldItalicFont=lmroman10-bolditalic.otf]
\setmonofont{lmmono10-regular.otf}[Scale=MatchLowercase]
\usepackage[margin=25mm]{geometry}
\usepackage[unicode,hidelinks]{hyperref}
\setlength{\emergencystretch}{3em}
\setlength{\parskip}{0.4em}
\setlength{\parindent}{0pt}
\allowdisplaybreaks[2]
\raggedbottom
\title{\texorpdfstring{The canonical vector Brownian innovations and initial Gaussian velocity}{The canonical vector Brownian innovations and initial Gaussian velocity}}
\author{LangevinLean mathematical exposition}
\date{}
\begin{document}\maketitle
\noindent Source module: \texttt{Kinetic\allowbreak{}Langevin/\allowbreak{}Canonical\allowbreak{}Vector\allowbreak{}Segments.\allowbreak{}lean}. The accompanying HTML pairs this mathematical exposition with the exact Lean statements and records the scope of the formalization.


The paper draws correlated velocity and position innovations in each spatial coordinate, independently across coordinates and time steps. Here every such draw is realized from the existing product of scalar canonical Brownian sources. The entire constructed array is proved to have exactly the previously defined Gaussian product law; equality of a few moments alone is not used as a substitute for equality of laws.


\section{\texorpdfstring{Transport independence along an actual law}{Transport independence along an actual law}}


A measurable map \(Y\) with actual law \(Q\) can pull back any independent measurable family defined on \(Q\). Lean proves this by the characterization of mutual independence through equality with the infinite product of marginal laws. It applies the pushforward composition theorem to the joint family and to each marginal. This elementary tool is used repeatedly to transfer established independence to the concrete Brownian source.



{\raggedright\noindent Lean declarations: \texttt{i\allowbreak{}Indep\allowbreak{}Fun\_\allowbreak{}comp\_\allowbreak{}has\allowbreak{}Law}.\par}

\section{\texorpdfstring{Build correlated scalar and spatial blocks}{Build correlated scalar and spatial blocks}}


A scalar block consists of all deterministic response kernels applied to the same Brownian interval. Packing their values with \texttt{With\allowbreak{}Lp.\allowbreak{}to\allowbreak{}Lp} gives a finite Euclidean vector whose law is the Gram Gaussian already proved in \texttt{Brownian\allowbreak{}Segments}. A vector block evaluates that scalar block on each spatial coordinate source. Every map is measurable because each L\ensuremath{{}^{2}} representative is strongly measurable, coordinate evaluation is measurable, and finite or infinite coordinatewise packing preserves measurability.



{\raggedright\noindent Lean declarations: \texttt{scalar\allowbreak{}Grid\allowbreak{}Block}, \texttt{measurable\_\allowbreak{}scalar\allowbreak{}Grid\allowbreak{}Block}, \texttt{vector\allowbreak{}Grid\allowbreak{}Block}, \texttt{measurable\_\allowbreak{}vector\allowbreak{}Grid\allowbreak{}Block}.\par}

\section{\texorpdfstring{Prove the law of the entire array}{Prove the law of the entire array}}


For one spatial source, disjoint Brownian intervals give independent scalar blocks. Different spatial sources are independent by the explicit product construction. The theorem \texttt{i\allowbreak{}Indep\allowbreak{}Fun\_\allowbreak{}uncurry'} combines these two facts into independence over iteration-coordinate pairs; reindexing swaps their order. Each marginal has the established Gram-Gaussian law, so \texttt{has\allowbreak{}Law\_\allowbreak{}infinite\allowbreak{}Pi} identifies the full array distribution. Currying the array is a measurable equivalence, and the product-measure curry theorem identifies the distribution as an independent stream of complete spatial blocks. Pullback along this actual path law proves independence across time steps.



{\raggedright\noindent Lean declarations: \texttt{vector\allowbreak{}Grid\allowbreak{}Block\_\allowbreak{}array\_\allowbreak{}independent}, \texttt{vector\allowbreak{}Grid\allowbreak{}Block\_\allowbreak{}array\_\allowbreak{}law}, \texttt{vector\allowbreak{}Grid\allowbreak{}Block\_\allowbreak{}path\_\allowbreak{}law}, \texttt{vector\allowbreak{}Grid\allowbreak{}Block\_\allowbreak{}time\_\allowbreak{}independent}.\par}

\section{\texorpdfstring{Identify the existing KLMC Gaussian stream exactly}{Identify the existing KLMC Gaussian stream exactly}}


For \(d\) finite spatial coordinates, the pair-indexed array has exactly the existing \texttt{kernel\allowbreak{}Noise\allowbreak{}Path} law. Specializing the response family to \((\psi_0,\psi_1)\) and packing each response into a vector produces the actual unscaled velocity-position innovation. Composition of the stream law with the existing \texttt{klmc\allowbreak{}Innovation} map proves its exact distribution. Independence of the full vector blocks survives this measurable packing. Each vector coordinate lies in L\ensuremath{{}^{2}} by the Brownian isometry; pulling it back along spatial evaluation and using \texttt{Mem\allowbreak{}Lp.\allowbreak{}of\_\allowbreak{}eval\_\allowbreak{}pi\allowbreak{}Lp} proves both phase components have finite second moments.



{\raggedright\noindent Lean declarations: \texttt{vector\allowbreak{}Grid\allowbreak{}Block\_\allowbreak{}kernel\allowbreak{}Noise\allowbreak{}Path\_\allowbreak{}law}, \texttt{canonical\allowbreak{}KLMCInnovation}, \texttt{measurable\_\allowbreak{}canonical\allowbreak{}KLMCInnovation}, \texttt{canonical\allowbreak{}KLMCInnovation\_\allowbreak{}law}, \texttt{canonical\allowbreak{}KLMCInnovation\_\allowbreak{}independent}, \texttt{canonical\allowbreak{}KLMCInnovation\_\allowbreak{}mem\allowbreak{}Lp}.\par}

\section{\texorpdfstring{Adjoin the independent initial Gaussian velocity}{Adjoin the independent initial Gaussian velocity}}


The source for a sampler is \(Q=\mathcal N(0,I_d)\otimes P_{\mathrm{Brownian}}\). Its first projection is the actual standard Gaussian initial velocity. Its second projection carries all constructed innovations. Product-measure projection laws prove the velocity law and preserve the innovation law. The product independence theorem proves that initial velocity is independent of the entire future innovation stream, while measure-preserving pullback gives all required L\ensuremath{{}^{2}} bounds. An arbitrary initial-position coupling can subsequently be extended by this same fixed source, independently of that coupling.



Pullback along the second projection also proves mutual independence of complete innovations across all iterations on this enlarged source; the new Gaussian velocity does not alter the noise law.



{\raggedright\noindent Lean declarations: \texttt{canonical\allowbreak{}Source\allowbreak{}Probability}, \texttt{canonical\allowbreak{}Source\allowbreak{}Probability\_\allowbreak{}is\allowbreak{}Probability\allowbreak{}Measure}, \texttt{canonical\allowbreak{}Source\allowbreak{}Velocity\_\allowbreak{}law}, \texttt{canonical\allowbreak{}Source\allowbreak{}Velocity\_\allowbreak{}mem\allowbreak{}Lp}, \texttt{canonical\allowbreak{}Source\allowbreak{}Innovation\_\allowbreak{}law}, \texttt{canonical\allowbreak{}Source\allowbreak{}Velocity\_\allowbreak{}innovations\_\allowbreak{}independent}, \texttt{canonical\allowbreak{}Source\allowbreak{}Innovation\_\allowbreak{}independent}, \texttt{canonical\allowbreak{}Source\allowbreak{}Innovation\_\allowbreak{}mem\allowbreak{}Lp}.\par}

\end{document}
